\section{Comparison for arbitrary cubic sides}\label{sec:different-sides}

The comparison states constructed above require increasingly divisible
periods. The original cubic sequence need not have this property.
We compare its side $n$ with a divisible side $m\le n$, using an
empty slab to absorb the difference in each of the three folding
directions. The empty slabs contribute partition functions without
insertions. Their logarithms are controlled by a common pressure
bound, and their total contribution is measured by their volumes
multiplied by the corresponding H\"older exponents. This is the
unequal-slab and companion-sequence argument of
\cite[Sections~11.1--11.5]{ChamberQED}, restricted to the original
cubic laws used in \cref{thm:main}.

Throughout this section the parameters retain their fixed values
$q>0$, $0<\rho<1$, $M\in\R$ and $w_i>0$. Write
$D_m=D(m/2,m,m)$ when $m$ is even. The original torus is always
$C(n)$; all tilted laws and comparison states in this section are
formed on the auxiliary tori $D_m$.

\subsection{Unequal slabs}\label{sec:unequal-heights}

Let $\ell(x)=n_\ell\cdot x$ be a mirror family and let
$\Xi=(\Gamma_\parallel,\chi_\parallel)$ be a cross-section. Put
\[
 |\Xi|=[\Z^3\cap\{\ell=0\}:\Gamma_\parallel].
\]
If $P>0$ belongs to $|n_\ell|^2\Z$, set
\[
 u_P=\frac{P}{|n_\ell|^2}n_\ell,\qquad
 \mathbb T(\Xi;P)=
   (\Gamma_\parallel\oplus\Z u_P,\chi_P),
 \qquad \chi_P|_{\Gamma_\parallel}=\chi_\parallel,
 \quad\chi_P(u_P)=-1.
\]
We use this notation only when the displayed period lattice has systole
at least four. Then $u_P$ is an integer vector, the torus is adapted
to $\ell$, and its volume is $P|\Xi|$. More generally, a torus with
this cross-section and height period $p$ has volume $p|\Xi|$.
Indeed, $\ell(\Z^3)=\Z$, and the quotient is counted by its $p$
height classes and the $|\Xi|$ classes in each height. The character
does not affect this count.

\begin{proposition}[Unequal slab folds]\label{prop:unequal-slabs}
Let $\mathbb T=(\Gamma,\chi)$ have cross-section $\Xi$ in the family
$\ell$, and let $g\in\Gamma$ be a closing vector of minimal positive
height $p=\ell(g)$. Choose an integer $k\ge1$ and admissible positions
\[
 c_0<c_1<\cdots<c_{2k}=c_0+p,\qquad h_j=c_j-c_{j-1}.
\]
Assume
\begin{equation}\label{eq:sides-systole}
 \min_{\gamma\in\Gamma_\parallel+
       \Z(2kh_j n_\ell/|n_\ell|^2),\ \gamma\ne0}
       |\gamma|_\infty\ge4,\qquad 1\le j\le2k.
\end{equation}
Let $G_j$ be a bounded scalar function of finitely many curvature
plaquettes contained in the open slab $c_{j-1}<\ell<c_j$.
For each $j$ put $\mathbb T_j=\mathbb T(\Xi;2kh_j)$ and form
$\operatorname{fold}_\ell G_j$ by successive reflections in the
planes $\ell=c_{j-1}+ih_j$. Then
\begin{align}
 Z_{\mathbb T_j}[\operatorname{fold}_\ell G_j]
   &=\|\mathfrak S[G_j]\|_{2k}^{2k}\ge0,
       \label{eq:sides-slab-norm}\\
 \left|Z_{\mathbb T}\left[\prod_{j=1}^{2k}G_j\right]\right|
   &\le\prod_{j=1}^{2k}
       Z_{\mathbb T_j}[\operatorname{fold}_\ell G_j]^{1/(2k)},
       \label{eq:sides-holder}\\
 |\mathbb T|&=\sum_{j=1}^{2k}\frac{|\mathbb T_j|}{2k}.
       \label{eq:sides-volumes}
\end{align}
In particular, the original torus need not admit a straightened torus
of height $p$.
\end{proposition}

\begin{proof}
Admissible positions form the group
$\frac12|n_\ell|^2\Z$. Thus $h_j$ belongs to this group,
$2kh_j\in|n_\ell|^2\Z$, and every plane used in the $j$th fold
is admissible. Condition \eqref{eq:sides-systole} therefore defines
each $\mathbb T_j$. Its volume is $2kh_j|\Xi|$, so summing
$h_j$ proves \eqref{eq:sides-volumes}.

The boundary spaces and the slab operator $\mathfrak S[G_j]$ depend
on the cross-section and the two endpoints, not on the torus used
to close a chain of slabs. Apply the reflected-pair identity
\eqref{eq:slab-fold-holder} on $\mathbb T_j$, whose normal character
is $-1$. Its closing translation cancels the translation of the
repeated pairs, giving
\[
 Z_{\mathbb T_j}[\operatorname{fold}_\ell G_j]
   =\operatorname{Tr}
       (\mathfrak S[G_j]^*\mathfrak S[G_j])^k.
\]
This proves \eqref{eq:sides-slab-norm}.

On the original torus the exact trace \eqref{eq:slab-trace} is
\[
 Z_{\mathbb T}\left[\prod_jG_j\right]
  =\operatorname{Tr}\left(
     V_g^{-1}U_{\mathrm f}^{(1+\chi(g))/2}
     \mathfrak S[G_{2k}]\cdots\mathfrak S[G_1]\right).
\]
Its closing factor is unitary. All slab operators are
Hilbert--Schmidt by \Cref{thm:slab-operators}, hence belong
to the Schatten class of exponent $2k$. Apply
\eqref{eq:slab-schatten-holder} with every exponent equal to $2k$,
and then use \eqref{eq:sides-slab-norm}. This gives
\eqref{eq:sides-holder}. The bounded scalar insertions are read on
periodically extended curvatures throughout; they require no choice
of a fermionic box or local lift.
\end{proof}

For example, when $n$ is odd the first diagonal height period of
$C(n)$ is $n$, whereas a straight normal closing vector would be
$(n/2)e_-$. This is not an integer vector. In the preceding proposition
each replacement normal period is $2kh_j$, which is even for integer
diagonal slab heights. This is the step that permits odd original sides.

\subsection{An empty slab and its pressure cost}\label{sec:empty-slab}

The exceptional torus in an unequal fold may be long. We use a bound
valid for every shape. Set
\begin{equation}\label{eq:sides-constants}
 H_M=|M|+\frac{1+\rho}{2}\sum_{i=0}^3w_i,
 \qquad \bar f=192\log H_M+3\log(r_qC_q).
\end{equation}
Both constants are finite, and $H_M>0$.

\begin{lemma}\label{lem:uniform-ceiling}
Every torus $\mathbb T$ used in the finite model satisfies
\[
 0<Z_{\mathbb T}[1]\le e^{\bar f|\mathbb T|}.
\]
Consequently $f_\infty\le f(a,b,c)\le\bar f$ for all integer
$a,b,c\ge4$.
\end{lemma}

\begin{proof}
The determinant bound is $\det\mathcal D_A^{\mathbb T}
\le H_M^{192|\mathbb T|}$, and the Maxwell integral is at most
$(r_qC_q)^{3|\mathbb T|}$. Together with the positive lower bound
in \eqref{eq:reg-partition-bounds}, these give the first assertion.
The second follows from the definition of pressure and
\Cref{prop:slab-pressure}.
\end{proof}

Use the diagonal cells \eqref{eq:slab-cells}, with heights
$(\ell^1,\ell^2,\ell^3)=(x_1-x_2,x_1+x_2,x_3)$, admissible
offsets $o_a$, and even width $h$. A finite family
$\mathfrak G=(G_\kappa)_{\kappa\in\mathcal K}$ will mean bounded
scalar curvature functions with $G_\kappa$ supported in the open
cell $Q_\kappa$. For a direction $a$ and a layer $j$, write
$\mathfrak R_a^{j,2k}\mathfrak G$ for the family obtained by retaining
the factors with $\kappa_a=j$ and reflecting that layer into the
$2k$ layers with indices $j,\ldots,j+2k-1$. The other two indices
are unchanged. Its product is exactly the slab fold of the original
layer product. This follows because the unique wall-group element
taking one layer to another acts on each of its cells by the
corresponding cell reflection.

\begin{lemma}[One empty slab]\label{lem:one-step}
Fix a direction $a$, a torus $\mathbb T$, and its cross-section $\Xi$
and height period $p$ for $\ell^a$. Suppose $(2k-1)h<p$, and put
\[
 h^{\mathrm{exc}}=p-(2k-1)h,\qquad
 \mathbb T^{\mathrm{reg}}=\mathbb T(\Xi;2kh),\qquad
 \mathbb T^{\mathrm{exc}}=\mathbb T(\Xi;2kh^{\mathrm{exc}}).
\]
Assume both replacement period lattices have systole at least four.
If every index of $\mathfrak G$ satisfies $0\le\kappa_a\le2k-2$,
then
\begin{align}
 \left|Z_{\mathbb T}[\prod\mathfrak G]\right|
 &\le Z_{\mathbb T^{\mathrm{exc}}}[1]^{1/(2k)}
    \prod_{j=0}^{2k-2}
      Z_{\mathbb T^{\mathrm{reg}}}
       [\prod\mathfrak R_a^{j,2k}\mathfrak G]^{1/(2k)},
       \label{eq:sides-step}\\
 |\mathbb T|
 &=\frac{2k-1}{2k}|\mathbb T^{\mathrm{reg}}|
     +\frac1{2k}|\mathbb T^{\mathrm{exc}}|.
       \label{eq:sides-step-volumes}
\end{align}
Every folded partition function on the right is nonnegative.
\end{lemma}

\begin{proof}
Cut at $o_a+jh$ for $0\le j\le2k-1$ and at $o_a+p$.
The first $2k-1$ slabs have height $h$; the final slab has height
$h^{\mathrm{exc}}>0$ and contains no factor. The positions are
admissible, since $p$ and $h$ are integers and $o_a$ is admissible.
Apply \Cref{prop:unequal-slabs} to their layer products
and to the constant $1$ in the final slab. The regular folds are
the products of $\mathfrak R_a^{j,2k}\mathfrak G$, while the empty
slab folds to $1$. Its three conclusions give the assertions.
\end{proof}

The empty slab has no local observable to transport. Thus its reflected
copies may have a different width, while every regular factor remains
in the original cell structure. Folding in direction $a$ leaves the
other two cell indices unchanged, so the same lemma can be applied
successively in the remaining directions.

\subsection{The comparison inequality}

Fix a real local trigonometric polynomial $O$ supported in a box
$y+[-r,r]^3\subset Q_0$, with $r\ge1$. Choose the cell width even,
$h\ge18$, and $h>8r+8$. The offsets are
$o_1=c_0,o_2=c_+,o_3=c_3$, with $c_0,c_+\in\Z$ and
$c_3\in\Z+1/2$. For $4h\mid m$, let $O_m^{\mathrm{fold}}$ be
the full reflected array on $D_m$, and let $\mu_{D_m}^t$ be the
tilt \eqref{eq:slab-tilted-law}. Its number of cells is
$N=m^3/|Q|$, where $|Q|=h^3/2$.

\begin{proposition}\label{prop:fold-any-side}
Suppose $4h\mid m$, $m\ge4h$, and $n\ge m$ is an integer.
Put
\begin{align*}
 V^{\mathrm{reg}}&=\frac{(m-h)^2(2m-h)}2,
 &V^{\mathrm{exc}}&=n^3-V^{\mathrm{reg}},\\
 \Delta'(n,m,h)&=f(m/2,2m,m)-f_\infty
       +\frac{V^{\mathrm{exc}}}{V^{\mathrm{reg}}}
                     (\bar f-f_\infty).
\end{align*}
Then $\Delta'(n,m,h)\ge0$ and, for every $t>0$,
\begin{equation}\label{eq:sides-expectation-comparison}
 \mu_{C(n)}(O)\le\mu_{D_m}^t(O)
                   +\frac{|Q|\Delta'(n,m,h)}t.
\end{equation}
There is no parity restriction on $n$.
\end{proposition}

\begin{proof}
Set $s=n-m\ge0$ and
\[
 k_1=k_3=\frac{m}{2h},\qquad k_2=\frac mh,\qquad
 \mathcal K^{\mathrm{reg}}=
   \{\kappa:0\le\kappa_a\le2k_a-2,\ a=1,2,3\}.
\]
These are integer ranges, with every $k_a\ge2$. Write
\[
 R=\prod_{a=1}^3(2k_a-1),\qquad
 N=8k_1k_2k_3=\frac{m^3}{|Q|}.
\]
Then $R|Q|=V^{\mathrm{reg}}$, and
$0<V^{\mathrm{reg}}<m^3\le n^3$. Since
$f(m/2,2m,m)\ge f_\infty$ and $\bar f\ge f_\infty$,
the claimed sign of $\Delta'$ follows.

\emph{The array on the original torus.}
For $\kappa\in\Z^3$ let
\[
 v(\kappa)=\frac h2(\kappa_1e_-+\kappa_2e_+)
                         +h\kappa_3e_3,\qquad
 G_\kappa=e^{tO^{\tau_{v(\kappa)}}}.
\]
The vector $v(\kappa)$ is integral because $h$ is even, and its
three height coordinates are $h\kappa$. Thus $G_\kappa$ is
supported in $Q_\kappa$. Define
\[
 \widetilde O_{\mathrm{reg}}
     =\sum_{\kappa\in\mathcal K^{\mathrm{reg}}}
                               O^{\tau_{v(\kappa)}},\qquad
 \mathfrak G=(G_\kappa)_{\kappa\in\mathcal K^{\mathrm{reg}}}.
\]
Then $\prod\mathfrak G=e^{t\widetilde O_{\mathrm{reg}}}$.
Jensen's inequality and translation invariance give
\begin{equation}\label{eq:sides-jensen}
 tR\mu_{C(n)}(O)
 \le\log Z_{C(n)}[e^{t\widetilde O_{\mathrm{reg}}}]
                            -\log Z_{C(n)}[1].
\end{equation}

\emph{The three replacement geometries.}
Apply \Cref{lem:one-step} in the order
$(\ell_-,\ell_+,x_3)$. The regular tori are
\[
 \mathbb T_0=C(n),\quad
 \mathbb T_1=D'(m/2,n,n),\quad
 \mathbb T_2=D(m/2,m,n),\quad
 \mathbb T_3=D_m.
\]
For the first step the cross-section of $C(n)$ is generated by
$ne_+,ne_3$, with characters $+1,-1$, respectively. Its height
period is $n$, with closing vector $ne_1$. The regular normal
vector is $k_1he_-=(m/2)e_-$, and the exceptional slab has height
$n-(2k_1-1)h=h+s$. Hence
\[
 \mathbb T_1^{\mathrm{exc}}=D'(k_1(h+s),n,n).
\]
This calculation uses only that $n$ is an integer.

In the second direction, $\mathbb T_1$ has cross-section generated
by $(m/2)e_-,ne_3$, both with character $-1$. Its height period is
$2n$, and the original closing vector $ne_+$ has character $+1$.
The trace therefore contains the parity operator, which is unitary
and is retained in \Cref{prop:unequal-slabs}. The regular
normal vector is $k_2he_+=me_+$, now with character $-1$.
The empty slab has height $h+2s$, giving
\[
 \mathbb T_2^{\mathrm{exc}}=D(m/2,k_2(h+2s),n).
\]
In the third direction, the cross-section of $\mathbb T_2$ is
generated by $(m/2)e_-,me_+$, both with character $-1$, and the
height period is $n$. The regular closing vector is $2k_3he_3=me_3$.
The empty slab has height $h+s$, so
\[
 \mathbb T_3^{\mathrm{exc}}=D(m/2,m,2k_3(h+s)).
\]

All these tori are admissible. For positive integers $a,b,c$,
the lattice generated by $ae_-,be_+,ce_3$ has systole
$\min(a,b,c)$: a vector has coordinates $(aA+bB,-aA+bB,cC)$,
and the maximum of the first two absolute values is
$|aA|+|bB|$. The stated minimum follows, and is attained by a
generator. Here every displayed coefficient is at least
$\min(m/2,2(h+s),n)\ge36$. This checks the regular and exceptional
systoles at each application of \Cref{lem:one-step}.

\emph{The reflected arrays.}
The first fold can be applied because the first indices of
$\mathfrak G$ lie in $[0,2k_1-2]$. After retaining a first layer
$j_1$ and reflecting it, the second and third index ranges are
unchanged. The second fold may therefore be applied to each
resulting family; afterwards the third range is still
$[0,2k_3-2]$, so the third fold applies as well.

For $j=(j_1,j_2,j_3)\in\mathcal K^{\mathrm{reg}}$, the final
family is the full reflected array of $G_j$ in the cells
\[
 j+\mathcal K_0,\qquad
 \mathcal K_0=\{\kappa:0\le\kappa_a\le2k_a-1\}.
\]
Indeed the three wall maps commute, and their product is the
unique wall-group element carrying $Q_j$ to the prescribed cell.
The periods of $D_m$ shift the cell indices by
$(2k_1,0,0)$, $(0,2k_2,0)$ and $(0,0,2k_3)$, so
$j+\mathcal K_0$ is a complete set of cell representatives.
Changing those representatives does not change a scalar fold on
periodic curvatures.

The full fold of $G_j$ has the same partition function as the full
fold of $e^{tO}$. Translation by $v(j)$ permutes the cell walls,
conjugates their reflected maps, and preserves $D_m$. Thus it
translates the full array of $e^{tO}$ into that of $G_j$, and
translation invariance of the finite functional gives
\[
 Z_{D_m}[\operatorname{Fold}_{D_m}G_j]
       =Z_{D_m}[e^{tO_m^{\mathrm{fold}}}].
\]

\emph{The weighted exceptional volume.}
The three H\"older steps give $R$ regular factors, each with exponent
$1/N$, and the exceptional exponents
\[
 \vartheta_1=\frac1{2k_1},\qquad
 \vartheta_2=\frac{2k_1-1}{4k_1k_2},\qquad
 \vartheta_3=\frac{(2k_1-1)(2k_2-1)}{8k_1k_2k_3}.
\]
All folded partition functions are nonnegative. Successive
substitution in \eqref{eq:sides-step} is therefore valid and yields
\begin{equation}\label{eq:sides-three-steps}
 Z_{C(n)}[e^{t\widetilde O_{\mathrm{reg}}}]
 \le\prod_{a=1}^3 Z_{\mathbb T_a^{\mathrm{exc}}}[1]^{\vartheta_a}
       Z_{D_m}[e^{tO_m^{\mathrm{fold}}}]^{R/N}.
\end{equation}
The three volume identities \eqref{eq:sides-step-volumes} are
\[
 |\mathbb T_{a-1}|=
       \frac{|\mathbb T_a^{\mathrm{exc}}|}{2k_a}
        +\frac{2k_a-1}{2k_a}|\mathbb T_a|,
        \qquad a=1,2,3.
\]
Substitute the third into the second, and then the result into the
first. Since $|\mathbb T_3|=m^3$ and $Rm^3/N=R|Q|$, this gives
\begin{equation}\label{eq:sides-exceptional}
 \sum_{a=1}^3\vartheta_a|\mathbb T_a^{\mathrm{exc}}|
       =n^3-R|Q|=V^{\mathrm{exc}}.
\end{equation}
Apply \Cref{lem:uniform-ceiling} separately to the exceptional
factors in \eqref{eq:sides-three-steps}. The resulting bound is
\begin{equation}\label{eq:sides-partition-comparison}
 Z_{C(n)}[e^{t\widetilde O_{\mathrm{reg}}}]
 \le e^{\bar f V^{\mathrm{exc}}}
       Z_{D_m}[e^{tO_m^{\mathrm{fold}}}]^{V^{\mathrm{reg}}/m^3}.
\end{equation}
The weighted identity \eqref{eq:sides-exceptional} is essential here.
The exceptional volumes themselves may exceed $n^3$.

\emph{Normalization and convexity.}
All partition functions in \eqref{eq:sides-partition-comparison}
are positive. For the inserted ones this also follows directly
from $|\widetilde O_{\mathrm{reg}}|\le R|O|_1$ and
$|O_m^{\mathrm{fold}}|\le N|O|_1$.
As in \eqref{eq:slab-fold-upper-tangent}, convexity of the logarithm
of the tilted partition function gives
\[
 \log Z_{D_m}[e^{tO_m^{\mathrm{fold}}}]
      \le\log Z_{D_m}[1]+tN\mu_{D_m}^t(O).
\]
The pressure bounds give
\[
 \log Z_{C(n)}[1]\ge n^3f_\infty,\qquad
 \log Z_{D_m}[1]\le m^3f(m/2,2m,m),
\]
where the second inequality is \eqref{eq:slab-partition-chain}
at the even side $m$. Combining these estimates with
\eqref{eq:sides-jensen} and
\eqref{eq:sides-partition-comparison} yields
\begin{align*}
 tR\mu_{C(n)}(O)
 &\le tR\mu_{D_m}^t(O)
       +V^{\mathrm{reg}}\bigl(f(m/2,2m,m)-f_\infty\bigr)\\
 &\hspace{35mm}+V^{\mathrm{exc}}(\bar f-f_\infty).
\end{align*}
Divide by $tR=tV^{\mathrm{reg}}/|Q|$ to obtain
\eqref{eq:sides-expectation-comparison}.
\end{proof}

\subsection{Divisible companion sides}\label{sec:companion-sequences}

The comparison side must accommodate every fixed cell width used in
the local estimates, including the coarser widths used for
amplification. Rounding to multiples of a single fixed integer would
not supply these widths. The following choice is eventually divisible
by every positive integer and changes the side by less than its square
root.

\begin{lemma}\label{lem:companions-exist}
Let $n_i\ge4$ be integers with $n_i\to\infty$. There are positive
integers $m_i$ such that $m_i\le n_i$ eventually, $m_i/n_i\to1$,
and every positive integer divides $m_i$ for all sufficiently large
$i$. One choice, for $n_i\ge64$, is
\[
 k(n)=\max\{k\in\Z:k\ge1,\ 8k!\le\sqrt n\},\qquad
 m_i=8k(n_i)!\left\lfloor\frac{n_i}{8k(n_i)!}\right\rfloor.
\]
The finitely many remaining terms may be set equal to $8$.
\end{lemma}

\begin{proof}
For $n\ge64$ the set defining $k(n)$ is nonempty and finite.
The remainder after division by $8k(n)!$ satisfies
\[
 0\le n-8k(n)!\left\lfloor\frac n{8k(n)!}\right\rfloor
       <8k(n)!\le\sqrt n.
\]
Thus $n_i-\sqrt{n_i}<m_i\le n_i$ for all large $i$, which gives
$m_i/n_i\to1$ and $m_i\to\infty$. Given an integer $d\ge1$,
the inequality $n_i\ge64(d!)^2$ holds eventually. It implies
$k(n_i)\ge d$, so $d$ divides $k(n_i)!$ and hence $m_i$.
\end{proof}

\begin{proposition}[Comparison for arbitrary original sides]
\label{thm:comparison-any-side}
Let $n_i\ge4$ tend to infinity, and choose $m_i$ as in
\Cref{lem:companions-exist}. Fix the cell data and observable
$O$ of \Cref{prop:fold-any-side}. Form the tilted laws
$\mu_i^t=\mu_{D_{m_i}}^t$ and their ordered comparison states using
these companion sides. Then
\[
 \Delta'_i=\Delta'(n_i,m_i,h)\longrightarrow0.
\]
If $\mu_{C(n_i)}\Rightarrow\nu$ on a prescribed subsequence and
$\nu_0$ is any comparison state subordinate to that subsequence,
then
\begin{equation}\label{eq:sides-limit-comparison}
 \nu(O)\le\nu_0(O).
\end{equation}
\end{proposition}

\begin{proof}
For all sufficiently large $i$, one has $4h\mid m_i$, $m_i\ge4h$
and $m_i\le n_i$. \Cref{prop:fold-any-side} therefore
applies. The pressure term $f(m_i/2,2m_i,m_i)-f_\infty$ tends
to zero by \Cref{prop:slab-pressure}. Moreover
\begin{equation}\label{eq:sides-volume-ratio}
 \frac{V_i^{\mathrm{reg}}}{n_i^3}
 =\left(\frac{m_i}{n_i}\right)^3
    \left(1-\frac h{m_i}\right)^2
    \left(1-\frac h{2m_i}\right)\longrightarrow1.
\end{equation}
Hence $V_i^{\mathrm{exc}}/V_i^{\mathrm{reg}}\to0$, proving
$\Delta'_i\to0$.

All divisibility requirements used to construct comparison states
hold for $(m_i)$. Choose the common volume subsequence $i_j$ and
strengths $t_k\downarrow0$ in \eqref{eq:diag-comparison-state},
subordinate to the prescribed original subsequence. For fixed $k$,
\eqref{eq:sides-expectation-comparison} gives
\[
 \mu_{C(n_{i_j})}(O)\le\mu_{i_j}^{t_k}(O)
                         +\frac{|Q|\Delta'_{i_j}}{t_k}.
\]
The observable $O$ is bounded and continuous. Letting $j\to\infty$
at fixed $t_k>0$ therefore gives $\nu(O)\le\nu_{t_k}(O)$.
Now let $k\to\infty$ to obtain \eqref{eq:sides-limit-comparison}.
The factor $1/t_k$ in the error is the reason for this order of limits.
\end{proof}

The law on the left of \eqref{eq:sides-limit-comparison} remains the
limit of the original cubic laws. The construction of $\nu_0$ uses
the companion volumes and retains the wall positivity and even
translation invariance proved for ordered comparison states.
